\chapter{User Stratification and Per-Group LoRA Adapter Routing}
\label{ch:groups}

\section{User Behavioral Clustering Objectives}

Chapter~\textbf{\ref{ch:time}} established that applying a globally fixed hourly loss weighting across all subscribers is suboptimal because temporal mobility distributions vary significantly between individuals. While certain subscribers commute during early morning hours, others move during late evening periods, and stationary subscribers exhibit minimal transition activity.

To resolve these temporal variations, two research questions were investigated:
1. \emph{Can subscribers be classified into distinct behavioral cohorts based on temporal transition profiles, using only historical prefix data available at inference time?}
2. \emph{Does conditioning model parameter optimization on identified user cohorts improve next-cell prediction accuracy over persistence baselines?}

\paragraph{Data Provenance.} Analysis evaluates dataset cohort \Dmob{} (\textbf{Table~\ref{tab:datasets}}), fitting behavioral profiles on \textbf{400 training subscribers} and evaluating performance on \textbf{100 held-out subscribers} across full-trajectory evaluation (persistence baseline score: 43.1\%).

\section{Feature Engineering for Temporal Mobility Profiles}

User stratification is established by constructing a four-dimensional temporal mobility vector for each subscriber. Physical base station mast transitions are computed across four daily intervals:
- Night (22:00--07:00)
- Morning (07:00--12:00)
- Afternoon (12:00--18:00)
- Evening (18:00--22:00)

Two key design constraints govern vector construction:
- \textbf{Mast-Level Transition Resolution}: Mobility is evaluated strictly at physical base station mast resolution, excluding intra-site sector handovers (as established in Chapter~\textbf{\ref{ch:data}}).
- \textbf{Inference-Time Causal Integrity}: Temporal mobility profiles are computed strictly from historical prefix segments available up to the prediction moment. Computing profiles on complete daily traces would introduce future information leakage.

\section{Taxonomy of Behavioral Groups and Per-Group Temporal Hardness}

Subscriber trajectories were partitioned into four behavioral clusters using $k$-means clustering evaluated on historical mobility vectors ($k=4$).

% F11 -- The four behavioural profiles, hour by hour.
\begin{figure}[!htbp]
\centering
\fitwidth{%
\begin{tikzpicture}
\begin{groupplot}[
  group style={group size=2 by 2, horizontal sep=1.5cm, vertical sep=1.95cm,
               xlabels at=edge bottom, ylabels at=edge left},
  rep, ybar, width=0.46\linewidth, height=3.9cm,
  /pgf/bar width=4.6pt, ymin=0, ymax=105, ytick={0,25,50,75},
  xmin=6.3, xmax=20.7, xtick={7,9,...,19},
  xlabel={hour of the day}, ylabel={changes of location per 100 records},
  y label style={font=\scriptsize},
  title style={font=\footnotesize,yshift=1pt},
  grid=major, ymajorgrids, xmajorgrids=false,
  enlarge x limits=false,
]
\nextgroupplot[title={\textcolor{clGzero}{$\blacksquare$} \textbf{G0} the stay-at-home user, 28\,\% of people}]
\addplot[fill=clGzero,draw=none] coordinates {
(7,33.6)(8,27.7)(9,29.2)(10,27.4)(11,28.0)(12,32.7)(13,31.3)(14,32.1)(15,31.3)(16,31.7)(17,35.0)(18,33.8)(19,29.3)(20,28.8)};
\node[font=\scriptsize,anchor=south,text=clGzero] at (axis cs:17,36) {35\,\%};
%
\nextgroupplot[title={\textcolor{clGone}{$\blacksquare$} \textbf{G1} the steady user, 35\,\%}]
\addplot[fill=clGone,draw=none] coordinates {
(7,48.4)(8,46.5)(9,44.1)(10,49.4)(11,47.2)(12,44.8)(13,42.3)(14,44.7)(15,42.1)(16,46.7)(17,48.7)(18,50.9)(19,46.4)(20,58.0)};
\node[font=\scriptsize,anchor=south east,text=clGone] at (axis cs:20,59) {58\,\%};
%
\nextgroupplot[title={\textcolor{clGtwo}{$\blacksquare$} \textbf{G2} the evening-moving user, 14\,\%}]
\addplot[fill=clGtwo,draw=none] coordinates {
(7,54.6)(8,57.9)(9,52.3)(10,39.9)(11,44.1)(12,53.4)(13,71.4)(14,73.1)(15,73.8)(16,79.7)(17,80.4)(18,81.9)(19,84.0)(20,79.2)};
\node[font=\scriptsize,anchor=south east,text=clGtwo] at (axis cs:19.6,85) {84\,\%};
%
\nextgroupplot[title={\textcolor{clGthree}{$\blacksquare$} \textbf{G3} the morning-moving user, 23\,\%}]
\addplot[fill=clGthree,draw=none] coordinates {
(7,77.1)(8,74.5)(9,71.1)(10,77.4)(11,79.0)(12,74.1)(13,71.4)(14,71.1)(15,73.0)(16,69.8)(17,72.0)(18,59.7)(19,40.6)(20,29.5)};
\node[font=\scriptsize,anchor=south,text=clGthree] at (axis cs:11,80) {79\,\%};
\end{groupplot}
\end{tikzpicture}}
\caption{The four kinds of user, hour by hour}
\label{fig:group-profiles}
\figtext{%
\figpara{What each bar is.}{The horizontal axis is the hour of the day. Each bar
counts, out of every 100 back-to-back records written during that hour, how many
landed on a \emph{different mast} from the one before: how often, at that hour,
these people actually went somewhere. The percentage in each title is the share
of \emph{people} in that kind, not of records. Nobody labelled the four kinds;
an automatic procedure was handed each person's rate over four broad stretches
of the day and asked for four average shapes. The names are mine, given
afterwards.}

\figpara{Why G2 and G3 are the whole point.}{G0 is flat and low, around 30
changes per 100 records and never above 35, so no hour of their day deserves
special treatment. G1 is flat and higher, about 45 all day. G2 climbs all
afternoon to 84 at seven in the evening; G3 does the exact opposite, 79 before
noon then down to 29. They are near mirror images, they hold 37\,\% of the
population between them, and \emph{any hour that is hard for one is easy for the
other}. One stretch of the day emphasised for everybody helps one of them and
misleads the other, which is exactly the failure of Chapter~\ref{ch:time}.
\smob, 400 training people; the four labelled peaks are the values printed on
the source chart and the remaining bars are recovered from its shape.}}

\end{figure}


\textbf{Figure~\ref{fig:group-profiles}} details the hourly base station transition rates (measured per 100 consecutive records) for each of the four identified behavioral groups ($G0$ to $G3$):

\begin{description}
  \item[G0 (Stationary Subscribers, 28\% of population).] Low, uniform transition activity throughout the day ($\approx 30$ transitions per 100 records in \textbf{Figure~\ref{fig:group-profiles}}, peaking at 35\% at 17:00).
  \item[G1 (Regular Mobile Subscribers, 35\% of population).] Moderate, uniform transition activity across daylight hours ($\approx 45$ transitions per 100 records in \textbf{Figure~\ref{fig:group-profiles}}, rising to 58\% at 20:00).
  \item[G2 (Evening-Mobile Subscribers, 14\% of population).] Low morning activity transitioning to peak mobility during late afternoon and evening hours ($\approx 84$ transitions per 100 records at 19:00 in \textbf{Figure~\ref{fig:group-profiles}}).
  \item[G3 (Morning-Mobile Subscribers, 23\% of population).] Peak mobility during early morning hours ($\approx 79$ transitions per 100 records before 12:00 in \textbf{Figure~\ref{fig:group-profiles}}), declining to 29.5\% during evening hours.
\end{description}

The inverse temporal mobility trends between $G2$ and $G3$ in \textbf{Figure~\ref{fig:group-profiles}} explain why globally fixed temporal loss weighting was ineffective: $G2$ and $G3$ represent near mirror images comprising 37\% of the total population, meaning any globally fixed temporal window that is hard for $G2$ is easy for $G3$.

% F12 -- The group label predicts difficulty, and it can be read from the past alone.
\begin{figure}[t]
\centering
\replegend{\swatch{clmove}~average site-change rate over the day \qquad
           \swatch{clbase}~persistence baseline \qquad
           \swatch{clmodel}~detection agreement}
\plotlink{\nbllm}{\fitwidth{%
\begin{tikzpicture}
\begin{groupplot}[
  group style={group size=2 by 1, horizontal sep=2.1cm},
  repbar, width=0.42\linewidth, height=5.4cm, xtick=data,
  title style={font=\footnotesize,yshift=1pt},
]
% ---- (a) mobility vs predictability ----
\nextgroupplot[title={(a) the ranking is exactly inverted},
  symbolic x coords={G0,G1,G2,G3},
  ymin=0, ymax=78, ytick={0,20,40,60},
  ylabel={\%}, /pgf/bar width=10pt]
\addplot[fill=clmove,draw=none] coordinates {(G0,33)(G1,42)(G2,51)(G3,52)};
\addplot[fill=clbase,draw=none] coordinates {(G0,56.3)(G1,41.7)(G2,19.7)(G3,17.8)};
%
% ---- (b) detection from the context prefix ----
\nextgroupplot[title={(b) group recovered from the past alone},
  symbolic x coords={ca,cb,cc,cd},
  xticklabels={30\%,50\%,80\%,90\%},
  xlabel={context fraction at question time},
  ymin=0, ymax=118, ytick={0,25,...,100},
  ylabel={agreement (\%)}, /pgf/bar width=13pt]
\addplot[fill=clnonsig,draw=none] coordinates {(ca,44)(cb,61)};
\addplot[fill=clmodel,draw=none] coordinates {(cc,92)(cd,92)};
\node[font=\scriptsize,text=clmodel,anchor=south,align=center] at (axis cs:cc,104) {reference protocol};
\end{groupplot}
\end{tikzpicture}}}
\caption{Group difficulty and group detectability}
\label{fig:group-difficulty}
\figtext{%
(a)~The group label is not merely descriptive: it says in advance how hard a
user will be. The orange bar is a group's average site-change rate over the
day and the grey-blue bar its persistence baseline, read against each other
rather than on a common scale. G0 changes site about a third of the time and
``repeat the last cell'' is already right 56.3\,\% of the time for them; G3
changes site half again as often and the same rule is right only 17.8\,\% of
the time. The ranking of persistence is exactly the inverse of the ranking of
mobility, which is worth verifying rather than assuming. (b)~At inference only
a user's past exists, so the group must be computed from the context prefix
and never from the target, or the whole idea is circular. The bars give how
often the group detected from the prefix alone agrees with the group computed
from the full sequence, at four context fractions. At the reference 80\,\%
context the pipeline recovers the right group for 92 of the 100 held-out users
unaided, and 90\,\% context gives the same figure; below roughly 60\,\% of the
day it collapses, 61\,\% agreement at half a day and 44\,\% at 30\,\%, which
is what the two pale bars mean. The mechanism is therefore worth using on a
user with most of a day behind them and worth nothing on a cold start, which
is exactly the case in which an operator has the least information and would
most want a prediction. \smob{}, 400 train / 100 test users.}
\figsource{\nbllm}{train\_cellid\_llm.ipynb}{the group-detection and per-group baseline cells}
\end{figure}


\textbf{Figure~\ref{fig:group-difficulty}} validates the efficacy of behavioral stratification. \textbf{\panelref{fig:group-difficulty}{a}} demonstrates an exact inverse relationship between physical group mobility (orange bars) and persistence baseline predictability (grey-blue bars): stationary $G0$ subscribers change location 33\% of the time and yield a 56.3\% baseline accuracy, whereas morning-mobile $G3$ subscribers change location 52\% of the time and yield a baseline accuracy of only 17.8\%. Overall, cluster assignments account for 72\% of predictability variance across held-out test subscribers. 

\textbf{\panelref{fig:group-difficulty}{b}} evaluates inference-time causal integrity by measuring how accurately user group assignments can be derived strictly from historical prefix segments. Given an 80\% prefix context, historical assignment matches full-day assignment for 92 out of 100 test subscribers (92\% accuracy). Below a 60\% context fraction, detection accuracy degrades rapidly (61\% at 50\% context, 44\% at 30\% context), indicating that early-stage trajectory routing requires alternative short-prefix classification methods.

\subsection{Data-Driven Derivation of Critical Mobility Windows}

Utilizing group-specific temporal profiles from \textbf{Figure~\ref{fig:group-profiles}}, critical transition windows were derived per group by identifying hourly intervals where the group's persistence baseline accuracy falls furthest below its daily average.

% T7 -- Per-group transition windows vs the single global guess (slides 51, 54).
\begin{table}[!htbp]
\centering\small
\caption{Hours derived per kind of user, and the hard moments they land on}
\label{tab:group-windows}
\begin{tabular}{@{}l>{\raggedright\arraybackslash}p{4.6cm}>{\raggedright\arraybackslash}p{6.1cm}@{}}
\toprule
\textbf{Kind of user} & \textbf{Hours derived from the data (cost multiplier)} & \textbf{What it says} \\
\midrule
G0 stay-at-home user & \textit{none} & a flat day has no hours worth emphasising \\
G1 steady user & 18--19h $\times1.4$, \ 20--21h $\times1.7$ & mild, and late in the day \\
G2 evening-moving user & 13--14h $\times1.7$, \ 16--21h $\times2.5$ & the five-hour block was later narrowed to 18--21h \\
G3 morning-moving user & 07--08h $\times2.0$, \ 10--13h $\times2.3$ & the mirror image of G2, and almost disjoint from it \\
\midrule
\multicolumn{3}{@{}l}{\textbf{Do those hours actually mark the hard moments?} Score of the one-line rule, inside versus outside}\\
\addlinespace
\multicolumn{2}{@{}l}{The guess of Ch.~\ref{ch:time}, 04--06h / 18--20h} & 43.6 inside vs 43.4 outside: a difference of \textbf{0.2}, so it \textbf{marks nothing} \\
\multicolumn{2}{@{}l}{The hours derived per kind of user} & 19.0 inside vs 46.4 outside: a difference of $\mathbf{27.4}$, \textbf{the genuinely hard hours} \\
\bottomrule
\end{tabular}

\end{table}


\textbf{Table~\ref{tab:group-windows}} lists the derived temporal transition windows and loss cost multipliers per group. For stationary $G0$ subscribers, no transition windows are upweighted since mobility remains uniformly low. For $G2$ and $G3$, derived transition windows are mutually disjoint, reflecting their opposing diurnal routines.

\textbf{Table~\ref{tab:group-windows}} confirms that group-derived temporal windows successfully isolate high-difficulty moments: baseline accuracy inside group-derived windows drops to \textbf{19.0\%} compared to \textbf{46.4\%} outside them, establishing a \textbf{27.4 percentage point separation}. In contrast, globally fixed heuristic windows (04:00--06:00 and 18:00--20:00) separated baseline difficulty by only 0.2 percentage points (43.6\% vs. 43.4\%), demonstrating that data-driven group stratification is necessary to identify true transition windows.

\section{Model Conditioning and Per-Group Adapter Architecture}

\subsection{Diagnostic Analysis of Textual Group Token Conditioning}

Initial conditioning experiments inserted explicit text tokens (\cid{<G0>} through \cid{<G3>}) at the beginning of input sequence prefixes. However, this textual conditioning model achieved 41.3\% accuracy against a 43.1\% persistence baseline (-1.8 percentage points).

Diagnostic ablation revealed the root cause: evaluating the model with correct group tokens (41.3\%), incorrect group tokens (41.2\%), or omitted group tokens (41.1\%) yielded virtually identical accuracies across a narrow 0.2 percentage point range. As illustrated in \textbf{\panelref{fig:adapter-routing}{a}}, self-attention layers diluted the single prefix group token across extended 500-token trajectory sequences, allowing the model to ignore the conditioning signal.

% F12b -- Group token vs. adapter routing: the change that worked.
\begin{figure}[!htbp]
\centering
\fitwidth{%
\begin{tikzpicture}[font=\small,
  seq/.style={draw=clrule,fill=clsoft,minimum height=0.55cm,inner xsep=4pt,rounded corners=1pt},
  gtok/.style={seq,draw=clGtwo,fill=clGtwo!30},
  base/.style={draw=clrule,fill=clnonsig!30,minimum width=3.2cm,minimum height=0.85cm},
  ad/.style={draw=clgain,fill=clgain!22,minimum width=1.25cm,minimum height=0.5cm,font=\scriptsize},
  adoff/.style={draw=clrule,fill=white,minimum width=1.25cm,minimum height=0.5cm,font=\scriptsize,text=clrule},
  lbl/.style={font=\scriptsize,text=clrule,align=center},
  ar/.style={-{Latex[length=4pt]},clrule},
]
% ================= (a) the group as a token: it failed =================
\node[lbl,anchor=south] at (2.4,2.35) {\textbf{(a) What failed:} the kind of user is one symbol in the text};
\node[gtok] (g) at (0,1.7) {\cid{<G2>}};
\node[seq,right=2pt of g] (s1) {\cid{<H07>}};
\node[seq,right=2pt of s1] (s2) {\cid{BSOCHO1}};
\node[seq,right=2pt of s2] (s3) {\dots};
\node[seq,right=2pt of s3] (s4) {\scriptsize $\dots$500 tokens};
\node[base] (b1) at (2.4,0.45) {one shared set of numbers};
\draw[ar] (2.4,1.4) -- (b1.north);
\node[lbl,anchor=north,text width=8.2cm] at (2.4,-0.15)
  {Correct kind 41.3\,\%, \emph{wrong} kind 41.2\,\%, \emph{no} kind 41.1\,\%.
   Three conditions within two tenths of each other: the model never read the
   symbol. One symbol among five hundred is free to ignore.};

% ================= (b) the group as a routing decision =================
\node[lbl,anchor=south] at (12.1,2.35) {\textbf{(b) What worked:} the kind of user chooses which numbers run};
\node[gtok] (g2) at (8.9,1.7) {G2};
\node[lbl,anchor=north,text=clGtwo] at (8.9,1.32) {identified from the past};
\node[seq,right=2pt of g2] (r1) {\cid{<H07>}};
\node[seq,right=2pt of r1] (r2) {\cid{BSOCHO1}};
\node[seq,right=2pt of r2] (r3) {\dots};

\node[base] (b2) at (12.1,0.75) {frozen base model};
\node[adoff] (a0) at (9.95,-0.35) {G0};
\node[adoff] (a1) at (11.35,-0.35) {G1};
\node[ad]    (a2) at (12.75,-0.35) {\textbf{G2}};
\node[adoff] (a3) at (14.15,-0.35) {G3};
\draw[ar] (b2.south) -- (12.1,0.05);
\draw[-{Latex[length=4pt]},clgain,thick] (g2.south) ++(0,-0.55) .. controls +(0,-1.0) and +(0,0.8) .. (a2.north);
\node[lbl,anchor=north,text width=8.2cm] at (12.1,-0.85)
  {Four independent sets of added numbers on the same frozen model. The kind of
   user is no longer text that can be diluted, it physically selects which
   numbers answer. Same people, same rule: from 1.8 behind to 3.7 ahead.};
\end{tikzpicture}}

\caption{The kind of user as a piece of text, and as a switch}
\label{fig:adapter-routing}
\figtext{%
\figpara{This is plumbing, not modelling, and it is what turned the project
around.}{The two panels are the same idea, delivered through two different
channels. Nothing else differs between them: same people, same protocol, same
model, same emphasised hours.}

\figpara{Panel (a): the kind of user as a piece of text.}{In the first attempt the
kind of user was written into the input as one extra symbol, \cid{<G0>} to
\cid{<G3>}, placed at the head of a sequence of up to five hundred symbols. The
model is free to look back at it from every later position, and it simply did
not. Feeding the correct kind, a deliberately wrong kind, or no kind at all gave
41.3, 41.2 and 41.1\,\%: three numbers inside two tenths of each other. Ignoring
that symbol cost the training almost nothing, so it was ignored.}

\figpara{Panel (b): the kind of user as a switch.}{In the second attempt the kind
of user is not information at all. Four separate sets of added numbers are trained
on the same frozen model, one per kind, and the kind identified from a person's
past decides which set runs. There is nothing left to dilute, because the
information never enters the text in the first place.}

\figpara{What that one change was worth.}{On the same people, the same protocol
and the same architecture, the distance to the one-line rule moved from 1.8
behind to 3.7 ahead (Figure~\ref{fig:group-gain}).}

\figpara{One caveat, stated plainly.}{Four sets of added numbers are also four
times the trainable capacity of one, so part of that move may be extra capacity
rather than the sorting itself. The control that would separate the two was not
run (Section~\ref{sec:retro-groups}).}}

\end{figure}


\textbf{Figure~\ref{fig:adapter-routing}} compares input prefix token conditioning in \textbf{\panelref{fig:adapter-routing}{a}} against physical parameter routing via separate LoRA adapters in \textbf{\panelref{fig:adapter-routing}{b}}.

\subsection{Implementation of Per-Group LoRA Adapter Routing}

To enforce group conditioning, per-group LoRA adapter routing was implemented (\textbf{\panelref{fig:adapter-routing}{b}}):
1. Base language model parameters remain frozen.
2. Four independent LoRA adapter modules are instantiated, corresponding to $G0$--$G3$.
3. Input trajectories are dynamically routed to the LoRA adapter corresponding to the subscriber's historical cluster assignment.

Physical parameter routing eliminates text token dilution by directing input trajectories to specialized subgroup weights, enabling collaborative parameter sharing among subscribers with identical mobility profiles.

\subsection{Performance Evaluation of Adapter Routing}

% F13 -- Where the gain lands once the group physically routes the sequence.
\begin{figure}[t]
\centering
\replegend{\swatch{clbase}~persistence baseline \qquad
           \swatch{clmodel}~model, group routed to its own adapter}
\plotlink{\nbllm}{\fitwidth{%
\begin{tikzpicture}
\begin{axis}[repbar, width=\linewidth, height=5.8cm,
  symbolic x coords={G0,G1,G2,G3,MEAN}, xtick=data,
  xticklabels={G0 homebody,G1 steady,G2 evening,G3 morning,mean of the four},
  /pgf/bar width=15pt,
  ymin=0, ymax=80, ytick={0,10,...,60},
  ylabel={top-1 accuracy (\%)},
  x tick label style={font=\small},
]
\addplot[fill=clbase,draw=none]  coordinates {(G0,56.0)(G1,42.0)(G2,19.7)(G3,17.8)(MEAN,33.9)};
\addplot[fill=clmodel,draw=none] coordinates {(G0,56.8)(G1,44.3)(G2,25.1)(G3,24.2)(MEAN,37.6)};
\node[font=\small,text=clgain,anchor=south] at (axis cs:G0,68) {$+0.8$};
\node[font=\small,text=clgain,anchor=south] at (axis cs:G1,56) {$+2.3$};
\node[font=\small,text=clgain,anchor=south] at (axis cs:G2,37) {$+5.4$};
\node[font=\small,text=clgain,anchor=south] at (axis cs:G3,36) {$+6.4$};
\node[font=\small,text=clgain,anchor=south] at (axis cs:MEAN,49.5) {$+3.7$};
\end{axis}
\end{tikzpicture}}}
\caption{Per-group gain after routing the group to its own adapter}
\label{fig:group-gain}
\figtext{%
The gain ranks in exactly the reverse order of the baseline. One pair of bars
per group, plus the unweighted mean of the four, with the persistence baseline
in grey-blue and the model in dark blue, scored on the same positions of the
same 100 held-out users over the full sequence. Nothing separates this run
from the one that scored $-1.8\pt$ except which weights were allowed to answer
(Figure~\ref{fig:adapter-routing}). G0, whom repetition already predicts
56.0\,\% of the time, gains $+0.8\pt$ and has almost nothing left to win; G3,
whom repetition predicts less than one time in five, gains $+6.4\pt$. The
ordering is exact, with no exception among the four, which is the claim; over
four points a correlation coefficient would add nothing to it. Read against
Chapter~\ref{ch:user} this is the
same finding arriving from the other direction: difficulty lives in the user,
and a model earns its cost precisely where the trivial rule fails. Two
qualifications belong with the headline. The $+3.7\pt$ is the unweighted mean
over four profiles counted once each; weighted by how many users each actually
holds it is $+3.3\pt$, which is the figure an operator would experience,
because 28\,\% of the population sits in the group that needed no model. And
the whole mechanism depends on the group being detectable, which
Figure~\ref{fig:group-difficulty}(b) shows collapsing below roughly 60\,\% of
a day of history. A third qualification, that four adapters are also four
times the trainable capacity and the control was never run, is in
Section~\ref{sec:retro-groups}.}
\figsource{\nbllm}{train\_cellid\_llm.ipynb}{the per-group adapter-routing run}
\end{figure}


Evaluating per-group LoRA adapter routing on active mobility test cohort \Dmob{} yields an overall top-1 accuracy of \textbf{46.8\%} against a persistence baseline of \textbf{43.1\%}, establishing a statistically significant \textbf{+3.7 percentage point gain} (unweighted average over the four groups; population-weighted gain is +3.3 percentage points, accounting for $G0$ representing 28\% of subscribers).

The five bar pairs in \textbf{Figure~\ref{fig:group-gain}} illustrate accuracy gains across the four user cohorts ($G0$--$G3$) and their overall average on the far right of \textbf{Figure~\ref{fig:group-gain}}. Accuracy improvements correlate inversely with baseline group predictability in \textbf{Figure~\ref{fig:group-gain}}: the largest gains manifest in high-mobility groups ($G3$ with +6.4 percentage points, 24.2\% vs. 17.8\%; $G2$ with +5.4 percentage points, 25.1\% vs. 19.7\%; $G1$ with +2.3 percentage points, 44.3\% vs. 42.0\%), whereas low-mobility groups ($G0$) perform closely to baseline bounds (+0.8 percentage points, 56.8\% vs. 56.0\%).

\section{Conclusion}
\label{sec:retro-groups}

Integrating user stratification with per-group LoRA adapter routing produced the first experimental model to surpass the persistence baseline across complete held-out evaluation trajectories (+3.7 percentage points in \textbf{Figure~\ref{fig:group-gain}}). Diagnostic analysis confirmed that routing trajectories to specialized subgroup parameters resolves text token dilution (\textbf{\panelref{fig:adapter-routing}{b}}) while enabling collaborative parameter sharing among subscriber cohorts.

